\settrack{tracktwo}
% VOICE: 1st-person personal (Bruce) — "I" not "Bruce". Personal/investigative voice.

\chapter{The Silence}
\label{pos29:the-silence}

\begin{quote}\small\textit{%
``Don't have good ideas if you aren't willing to be responsible for them.''%
}\par\raggedleft --- Perlis, Epigram \#95 (1982)\end{quote}

\vspace{0.5cm}

\srcnote{Session 20 UQ extraction (38 prompts) | session20-uq-extraction.md | 2026-02-24 | The Silence --- separation, constraint, break, [redacted], Europe, research}

\section*{The Separation}
\label{pos29:the-separation}

In September 2006, Healer moved out of the garage room we'd built him at the house on Fifteenth and Gant. We'd bought the construction materials together in 2004 and framed the room out of what had been storage from a previous life. For two years it was his space --- a sleeping bag, a laptop I'd bought him, and whatever he was working on that I didn't ask about.

The goodbye was fond but he didn't drag it out. Healer was not a man who lingered. He took a backpack of personal effects, the laptop, and the car I'd bought him. Zero money. I'd disbursed small amounts of cash over the years --- a few hundred here and there, as needed. Healer once told me that both his sister and his sensei had warned him against having a lot of money. When I knew him he behaved as if he had taken some sort of vow of poverty.

He said goodbye and he left. His decision, not an eviction. We were not on great terms but we were not on bad ones either. I felt bad. I silently wished him the best.

[REDACTED] In August 2006 I'd informed Healer that I could no longer assist him \& that I was impatient. [REDACTED]


I spoke with him one more time a bit later. He offered me the chance to help with a project and I declined. I regret it to this day.

I left Symantec in 2007 with severance. Under Possibility~C, I was never unsupported during the silent years --- I just didn't always know who was supporting me.

\section*{The Quiet Years}
\label{pos29:the-quiet-years}

I settled into a quieter life in Corvallis. The surface looked full. Underneath, it was constrained.

I could think about everything that had happened with Healer and the project. But I couldn't speak about it. The most important experience of my life, and no one to tell. No one around me noticed, because no one knew there was anything to notice.

Exile from yourself. That's the closest I can get to describing it.

\section*{The Loophole}
\label{pos29:the-loophole}

Dave Bannon and I had been fast friends since 2005 --- instant connection, the kind where you skip the small talk and go straight to the hard questions.

In 2008, Bannon started asking questions. Direct, pedagogical questions --- the kind a physicist asks when he's testing a hypothesis. He got a partial version of the story that year. He listened as a skeptical scientist, and was surprised that I had answers to all his questions --- answers he could independently verify. Of perhaps fifty people I told portions of this story to over twenty years, none finished the entire thing. The pedagogical prerequisites are too demanding. Dave came closest to full understanding, and even he needed Kauffman as a running start.

Bannon had read Kauffman independently. That was critical --- it meant he already had the conceptual vocabulary. He'd also arrived at his own independent intuition: that Google Search's performance was ``too good and too fast'' for classical hardware. A performance-based observation, not an architectural theory. He'd reached this conclusion before hearing anything from me.

Under Possibility~A, Bannon is a skeptical physicist who listened carefully and reserved judgment --- which is what good scientists do. Under Possibility~C, that reservation is itself significant.

From decades of trying to tell this story in person, three problems always intervene. First, knowledge prerequisites --- no one knows quantum physics, topological neural networks, SAS operations, and tradecraft. The story requires background in all four. Second, time --- the telling invariably stretches to many hours and people lose patience. Third, weirdness --- the concept is too strange for most people to process. Every attempt to share what I knew ran into this wall.

\section*{The Recognition}
\label{pos29:the-recognition}

Once the block broke, connections formed fast.

I became convinced that Healer's project became something I could see in the news --- something large and consequential that matched what he had described in vague terms. Someone else had done the job I'd been vetoed from doing. This is my interpretation, not something Healer told me. The specifics belong to the redacted chapter.

What I remember clearly is sitting in my living room in South Corvallis and realizing: I had information in my brain that the President of the United States did not. That thought made me nervous in a way that nothing else in this story had.

Under Possibility~A, a man with an overactive pattern-matching faculty saw a vague invitation and a public event and connected them after the fact. Under Possibility~B, Healer may have been peripherally connected to the events I later identified, but the connection is inflated. Under Possibility~C, the connection is exactly what it appears to be. The redacted chapter, when it is eventually published, will examine this claim in detail. Under Possibility~A, this redaction is theater --- there is nothing to redact. Under Possibility~B, there may be something, but the claimed sensitivity is overstated. Under Possibility~C, the author's caution is well-founded.

\section*{Europe}
\label{pos29:europe}

By 2010 there were reasons to be scared that I cannot fully explain in this edition. [REDACTED] I spent six months in Europe, working as a CTO.

\section*{Thread-Pulling}
\label{pos29:thread-pulling}

Early 2011. I returned from Europe and moved back to a community in South Corvallis. I didn't work except light farm labor. Instead I pulled threads.

Everything I'd learned from Healer --- the five sciences, the architecture, the pedagogy --- I began reconstructing from scratch. Relearning what I could, verifying what I could verify, mapping the public-domain trail.

At the Oregon Country Fair that summer I met Mark R, author of oilempire.us. Random encounter at Chew Rays cafe --- Teri and I were in the crowd while Mark was speaking. He saw me in the back and somehow knew instantly that I understood his content. Body language, maybe. I could have given the same lecture. One incisive question after. We exchanged contact information. Became friends.

Excepting Healer, Mark was the smartest person I knew. He had an intellectual edge on me in non-science areas. His specialization was the study of the Holocaust --- all his extended family had been killed in the Holodomor and then the Holocaust. Two full shelves of books on the subject. Kind, gentle, completely nonviolent.

I presented the TQNN theory to Mark as A/B/C after a few months. He put substantial time into his own research. Concluded: insufficient evidence to determine which possibility was correct, but plausible. The science was beyond him.

Then Mark searched the names I'd given him --- Kauffman, Wolfram --- and found Bill Joy's ``Why the Future Doesn't Need Us'' in \textit{Wired}. He sent me the link. This was 2012.

I recognized the article instantly. It was the fulfillment of something Healer had told me: that they had already told the world.

\section*{The Research Begins}
\label{pos29:the-research-begins}

I wrote tens of thousands of words about this. I began keeping careful notes in 2011, after I returned from six months in Europe working on a startup company.

In 2011 I wrote down everything I could remember in a single sitting. My intent was to document my hypothesis.

March 17, 2012: I posted to Cryptome. June 19, 2012: a comment on Slashdot under my real user ID\@. I built a toy neural network in JavaScript for biometric identification and wrote about it publicly --- for Slashdot, for Bruce Schneier's blog. Someone put the Cryptome post on a blog; my friend in Luxembourg saw it from across the Atlantic.

\vspace{0.5cm}

Under any possibility, this is a man who lost years to a story he cannot prove. Whether the story is true changes what we call it. Under Possibility~C, it is sacrifice. Under Possibility~A, it is tragedy of a different kind. Under Possibility~B, it is something in between --- a man who touched something real and spent years trying to name it.

\vspace{1cm}\noindent\rule{0.5\textwidth}{0.5pt}\vspace{0.5cm}

\begin{center}\textit{This story will be my life's work.}\end{center}

\chapterreturn
